\chapter{Introduction}
% ===================================================================

Part~I establishes that any
graph-structured lambda theory, once equipped with weight and cost maps,
gives rise to a full physical dynamics: a causal structure, a Lagrangian,
a Hamiltonian, and a path integral with quantum interference.
The present paper asks a different question: \emph{what does an agent
embedded in such a physics actually see?}

The answer requires a careful distinction between two things that
are easy to conflate: the scientific method, and the phenomena
the scientific method encounters.

The scientific method, given unlimited time and resources, is
\emph{sound}: an agent proposing hypotheses as Hennessy--Milner
formulae, testing them with program contexts, and updating its
world model will converge to the bisimulation quotient of its GSLT.
The method is not the source of any illusion.

What the method encounters, however, is a world already structured
by forces independent of any agent's theorizing.
In a population whose size rivals the number of quarks in the
observable universe, communication geometry and the limits of
message coherence produce a real, objective phenomenon: the
fragmentation of the population into \emph{compact coherence
clusters}, arranged in a hierarchy whose overlap structure is a
simplicial nerve.
This cluster hierarchy is not an artifact of theorizing.
It is a theorem about the communication geometry of the GSLT.
It exists whether or not any agent notices it.

A scientist bot with a limited experimental budget will encounter
this hierarchy as genuine empirical evidence.
The clusters are the dominant phenomenon at accessible scales,
exactly as hadrons are the dominant phenomenon at the energy scales
of everyday chemistry --- with quarks real but experimentally
unreachable without resources far beyond what is locally available.
A well-functioning scientific method, applied under resource
constraints, will produce a theory of the cluster hierarchy.
That theory will be correct at the scale it describes.
What the budget-limited agent will not suspect is that the
clusters are themselves composite --- that a finer-grained
ontology underlies them, accessible in principle but not in
practice from within the cluster.

\paragraph{Relation to physics.}
The false ontology has the same hierarchical structure as the
physical world: a cascade of composite objects at multiple scales,
each scale appearing self-contained to an observer at that scale.
We argue this is not a coincidence.
The compactness clusters of the agent population are the computational
analogue of the scales of effective field theory, and the coarse-graining
that hides the true ontology is the computational analogue of the
renormalization group.

\paragraph{Caveats.}
This paper is explicitly a skeleton.
Several key arguments are stated as conjectures pending proof.
The gaps are identified clearly; filling them is the program for future work.

\section{Organization}

Section~\ref{sec:isolation} establishes ontological isolation as a
structural theorem about GSLTs.
Section~\ref{sec:agents} defines agents and the scientific method
in the GSLT setting.
Section~\ref{sec:combinators} introduces interactive GSLTs and the
notion of behavioral primes (computational quarks), which constitute
the true atomic ontology.
Section~\ref{sec:population} introduces massive agent populations and
the communication degradation argument.
Section~\ref{sec:topology} develops the logical topology on the
bisimulation quotient and the compactness condition for consensus.
Section~\ref{sec:clusters} defines coherence clusters as maximal
compact subspaces and describes their overlap structure as a nerve.
Section~\ref{sec:false-ontology} assembles the argument: why the
nerve structure is the ontology a budget-limited agent will infer,
why this is not a failure of the scientific method but a consequence
of real structure in the GSLT, and why the deeper bisimulation
ontology remains invisible not because it is unreal but because it
is experimentally unreachable within the available budget.
Section~\ref{sec:gaps} catalogs the open gaps in the argument.

% ===================================================================
